% !TeX root = ../../tesis.tex
%%%%%%%%%%%%%%%%%%%%%%%%%%%%%%%%%%%%%%%%%%%%%%%%%%%%%%%%%%%%%%%%%%%
% SECCIÓN 5.4.1 — COEFICIENTE DE FRICCIÓN VIAL (CF)
%%%%%%%%%%%%%%%%%%%%%%%%%%%%%%%%%%%%%%%%%%%%%%%%%%%%%%%%%%%%%%%%%%%
\subsection{Coeficiente de Fricción Vial (\texorpdfstring{$CF$}{CF})}
\label{subsec:fase2-cf}

El Coeficiente de Fricción Vial, definido en §\ref{subsec:friccion-vial}
e implementado como \texttt{FrictionCalculator} en §\ref{subsec:vft-cf},
traduce el tipo de derecho de vía de cada segmento en un multiplicador
sobre el tiempo de traslado:

\begin{equation*}
    CF = 1 + \alpha \cdot \beta_{\text{ZMVM}}
    \tag{\ref{eq:vft-cf}}
\end{equation*}

Con $\beta_{\text{ZMVM}} = 0.759$ \citep{tomtom2024} y cuatro valores
de $\alpha$ calibrados según el grado de segregación física de la
infraestructura \citep{sedatu2019manual}, los coeficientes resultantes
para cada categoría son los del Cuadro~\ref{tab:cf-derecho-via}.

\begin{table}[H]
    \centering
    \caption[Coeficiente de Fricción Vial por categoría de derecho de vía]{%
    Coeficiente de Fricción Vial ($CF$) por categoría de derecho de vía
    y sistemas de transporte asignados en la red de la \zmvm{}.
    Un $CF = 1.000$ indica operación sin interferencia del tráfico.
    Corrida: 2026-04-18.}
    \label{tab:cf-derecho-via}
    \small
    \begin{tabular}{|l|c|r|l|}
        \hline
        \textbf{Derecho de vía} & \textbf{$\alpha$} &
        \textbf{$CF$} & \textbf{Sistemas} \\
        \hline
        Exclusivo  & 0.0 & 1.000 & Metro, Metrobús, Tren Ligero, Tren Suburbano \\
        Confinado  & 0.2 & 1.152 & Cablebús, Mexibús, Mexicable \\
        Compartido & 0.5 & 1.380 & Trolebús \\
        Mixto      & 1.0 & 1.759 & RTP, Corredores Concesionados, Pumabús \\
        \hline
    \end{tabular}
    \fuentefigura{Elaboración propia con \texttt{VFTModel}
    \citep{cabrera2026VFTModel};
    $\beta$: \citep{tomtom2024};
    $\alpha$: \citep{sedatu2019manual}.}
\end{table}

La brecha entre extremos es considerable: los sistemas con derecho de
vía exclusivo operan a su velocidad de diseño sin penalización, mientras
que RTP y los Corredores Concesionados —que concentran más del
85\,\% de las aristas de servicio real de la red— absorben el índice
de congestión completo de la ciudad. Dado que esas líneas son también
las que conectan las periferias con el centro, el $CF$ amplifica las
asimetrías de cobertura y capilaridad ya documentadas en la Fase~1.
